\clearpage
\onecolumn
\thispagestyle{empty}
\vspace*{\fill}
\begin{center}
	{\color{radblue}\rule{\linewidth}{2pt}}\par
	\vspace{6mm}
	{\Huge\bfseries\color{radblue} Abdomen}\par
	\vspace{6mm}
	{\color{radblue}\rule{\linewidth}{2pt}}
\end{center}
\vspace*{\fill}
\clearpage
\twocolumn
\pagestyle{contenido}


%%%%%%%%%%%%%%%%%%%%%%%%%%%%%%%%%%%%%%%%%%%%%%%%%%%%%%%%%%%
%  SECCIÓN: Abdomen
%  Archivo: secciones/abdomen.tex
%  Sin preámbulo ni \begin{document} — lo maneja main.tex
%%%%%%%%%%%%%%%%%%%%%%%%%%%%%%%%%%%%%%%%%%%%%%%%%%%%%%%%%%%
	
	\section{Abdomen}
	\subsection{Generalidades de Abdomen}
	
	%----------------------------------------------------------
	%  INFORMACIÓN COMPLEMENTARIA GENERAL
	%----------------------------------------------------------
	
	\subsubsection*{División por Cuadrantes}
	
	Dos planos perpendiculares trazados a nivel del ombligo dividen la región en nueve zonas:
	
	\begin{itemize}
		\item \textbf{Hipocondrio derecho:} hígado y vías biliares
		\item \textbf{Epigastrio:} estómago
		\item \textbf{Hipocondrio izquierdo:} bazo
		\item \textbf{Flanco derecho:} colon ascendente
		\item \textbf{Mesogastrio / región umbilical:} intestino delgado
		\item \textbf{Flanco izquierdo:} colon descendente
		\item \textbf{Fosa ilíaca derecha:} ciego y apéndice
		\item \textbf{Hipogastrio:} vejiga urinaria (cuando está llena)
		\item \textbf{Fosa ilíaca izquierda:} colon sigmoideo
	\end{itemize}
	
	\subsubsection*{Repères Anatómicos}
	
	\textbf{Límite superior:} cúpulas diafragmáticas. \\
	\textbf{Límite inferior:} sínfisis del pubis.
	
	\begin{itemize}
		\item \textbf{Apófisis xifoides:} T9--T10
		\item \textbf{Borde costal inferior:} L2--L3
		\item \textbf{Cresta ilíaca:} L4--L5
		\item \textbf{EIAS}
		\item \textbf{Trocánter mayor:} 3--4 cm por encima de la sínfisis púbica
		\item \textbf{Sínfisis del pubis:} unión anterior de ambos huesos pelvianos
		\item \textbf{Tuberosidad isquiática:} 1--4 cm por debajo de la sínfisis del pubis
	\end{itemize}
	
	\subsubsection*{Consideraciones Técnicas Generales}
	
	\textbf{Bajo contraste intrínseco:} estructuras de densidad similar; se usan técnicas más contrastadas (básica o intermedia).
	
	\textbf{Densidades presentes:}
	\begin{itemize}
		\item Aire (tubo digestivo)
		\item Grasa (rodea riñones, músculos psoas)
		\item Agua (vísceras)
		\item Hueso
	\end{itemize}
	
	\textbf{Líneas grasas:} interfase entre densidad de agua y grasa. Su visualización confirma una técnica adecuada. Principales referencias: borde ínfero-externo del hígado, línea de los músculos psoas (desde L1 hasta la pelvis), contorno perirrenal. Si las vísceras no se ven es porque no hay contraste.
	
	\textbf{Movimientos involuntarios:} inevitables; se compensan con tiempos de exposición cortos. kVp no muy altos, especialmente en pacientes con más grasa abdominal (genera radiación dispersa). El paciente debe estar relajado y con la respiración suspendida en espiración: disminuye el aire y el espesor.
	
	\subsubsection*{Preparación del Paciente}
	
	Indispensable para estudios programados. Dieta controlada, laxantes y enemas previos. Elimina gases y materia fecal del tubo digestivo, mejorando la calidad de imagen.
	
	\textbf{Sin preparación:} pacientes con abdomen agudo, sospecha de perforación o provenientes de emergencias.
	
	\subsubsection*{Consideraciones sobre los Enfoques}
	
	La imagen debe abarcar toda la región: desde cúpulas diafragmáticas hasta sínfisis del pubis. Si no es posible en una sola proyección, se realizan dos: una en decúbito supino de sínfisis hacia arriba y otra en bipedestación de cúpulas hacia abajo.
	
	No confundir decúbito lateral (para obtener un frente de abdomen con RC horizontal) con un paciente en decúbito lateral para un perfil.
	
	\subsubsection*{Protocolo Abdomen Agudo}
	
	\textbf{Par estándar:}
	\begin{enumerate}
		\item Frente en decúbito supino
		\item Frente en bipedestación (ortostática)
	\end{enumerate}
	
	Si el paciente no puede pararse: decúbito lateral izquierdo con rayo horizontal, o como complemento un frente de tórax (sentado). El decúbito lateral izquierdo sustituye la bipedestación en pacientes que no pueden deambular; es el más importante para detectar neumoperitoneo.
	
	\textbf{RC en bipedestación:} perpendicular, para no distorsionar los niveles hidroaéreos.
	
	\subsubsection*{Indicaciones por Hallazgo Clínico}
	
	\textbf{Cuerpo extraño:} frente y perfil de abdomen.
	
	\textbf{Niveles hidroaéreos:} de pie (preferible) o decúbito supino frente y perfil. En el perfil el líquido baja y el aire sube.
	
	\textbf{Neumoperitoneo:} complementar con frente de tórax. Paciente de pie o sentado con torso erguido. El aire tiende a ascender a la zona más alta.
	
	\textbf{Abdomen agudo:} AP decúbito supino + AP bipedestación.
	
	\textbf{Trauma abdominal - hallazgos a valorar:}
	\begin{itemize}
		\item Neumoperitoneo: indica perforación de víscera hueca
		\item Borramiento de línea del psoas: asociado a lesiones retroperitoneales
		\item Fracturas costales
	\end{itemize}
	
	%----------------------------------------------------------
	%  Los enfoques se agregan a continuación como \subsection{}
	%----------------------------------------------------------
	
	\clearpage
	
%%%%%%%%%%%%%%%%%%%%%%%%%%%%%%%%%%%%%%%%%%%%%%%%%%%%%%%%%%%
%  SECCIÓN: Abdomen
%  Archivo: secciones/abdomen.tex
%  Sin preámbulo ni \begin{document} — lo maneja main.tex
	%%%%%%%%%%%%%%%%%%%%%%%%%%%%%%%%%%%%%%%%%%%%%%%%%%%%%%%%%%%
	
	\section{Abdomen}
	
	\subsection{Abdomen AP en Decúbito Supino}
	
	\subsubsection*{1. Identificación}
	
	\textbf{Tipo:} Proyección básica. Forma par con la proyección AP en bipedestación.
	
	\textbf{Indicaciones clínicas principales:}
	\begin{itemize}
		\item Obstrucción intestinal y evaluación de niveles hidroaéreos en paciente sin capacidad de bipedestación
		\item Calcificaciones abdominales, ascitis y exploraciones con contraste
		\item Abdomen agudo
	\end{itemize}
	
	\subsubsection*{2. Técnica Base}
	
	\textbf{Receptor de imagen:} 35$\times$43 cm. Orientación \textbf{longitudinal} en todos los biotipos.
	
	\textbf{DFRI:} \textit{[dato no suministrado]}.
	
	\textbf{Técnica:} Intermedia. Las densidades del abdomen son muy homogéneas (bajo contraste intrínseco), por lo que se requiere una escala larga de grises con kVp intermedio para diferenciar el contorno de las vísceras (riñones), tejido adiposo, aire y hueso. En paciente obeso se utiliza técnica básica por el gran espesor y la consecuente mayor radiación dispersa que engrisa la imagen; en paciente pediátrico muy delgado puede homogeneizarse con herramientas de posproceso.
	
	\textbf{Factores técnicos:} kVp \textbf{[valor]} --- mAs \textbf{[valor]}. Medición del espesor a nivel de las crestas iliacas.
	
	\textbf{Bucky:} Bucky de mesa con grilla.
	
	\subsubsection*{3. Posicionamiento}
	
	\paragraph{Paciente}
	Decúbito supino, brazos a los lados del cuerpo, piernas extendidas con soporte bajo las rodillas. Plano medio sagital perpendicular y centrado al receptor. Se puede aplicar tracción desde los tobillos (verificar ausencia de prótesis) o desde los hombros para corregir la alineación vertebral.
	
		
	\begin{figure}[h]
		\centering
		\includegraphics[width=0.7\linewidth]{imagenes/abd2}
		\caption{}
		\label{fig:abd}
	\end{figure}
	
	
	\paragraph{Región de interés}
	Pelvis y abdomen sin rotación: EIAS equidistantes al receptor. El colimador se abre en su totalidad a lo ancho para incluir toda la región abdominal.
	
	\paragraph{Rayo Central}
	Perpendicular al plano del receptor. Punto de entrada: \textbf{a la altura de las crestas iliacas}.
	
	\paragraph{Receptor de imagen}
	Centrado a la altura de las crestas iliacas. El borde inferior debe incluir la sínfisis del pubis; el borde superior debe alcanzar las cúpulas diafragmáticas. Si el paciente es muy alto y no se logra abarcar toda la pelvis, se realiza una segunda exposición centrada en pelvis para incluir la vejiga.
	
	\paragraph{Respiración}
	Al final de la espiración. El descenso del diafragma permite visualizar las estructuras supradiafragmáticas y amplía el campo abdominal visible.

	
	\subsubsection*{4. Anatomía Visible}
	
	Contorno de riñones, borde inferior hepático, músculo psoas bilateral, marco intestinal, columna lumbar y pelvis ósea. Límite superior obligatorio: cúpulas diafragmáticas. Límite inferior obligatorio: sínfisis del pubis. Si alguno de los dos límites queda excluido, la radiografía carece de valor diagnóstico. En decúbito supino los niveles hidroaéreos pueden observarse pero no son tan evidentes como en bipedestación.
	
	\subsubsection*{5. Criterios de Evaluación}
	
	\textbf{Posicionamiento correcto:} Columna vertebral centrada en el receptor; crestas iliacas y cúpulas diafragmáticas incluidas; sínfisis del pubis visible en el borde inferior.
	
	\textbf{Ausencia de rotación:} Espinas iliacas anterosuperiores simétricas; alas iliacas y agujeros obturadores simétricos; apófisis espinosas centradas en los cuerpos vertebrales lumbares; caderas equidistantes a los bordes laterales de la imagen.
	
	\textbf{Penetración y contraste:} Escala larga de grises con kVp intermedio; diferenciación visible entre contorno renal, tejido adiposo perivisceral, gas intestinal y estructuras óseas.
	
	\subsubsection*{6. Puntos Críticos}
	
	\textbf{No incluir la sínfisis del pubis:} La vejiga y el segmento pélvico quedan fuera del campo; ante paciente de gran talla agregar una segunda proyección centrada en pelvis.
	
	\textbf{Rotación de pelvis u hombros:} Genera asimetría de alas iliacas y agujeros obturadores; corregir distribución del peso y alineación del plano sagital antes de la exposición.
	
	\textbf{Adaptaciones:} Si el paciente no puede mantenerse en bipedestación para evaluar aire libre, realizar el enfoque en decúbito lateral izquierdo con RC horizontal como alternativa.
	
	\textbf{Trucos:} La tracción desde los tobillos mejora la alineación vertebral; si hay prótesis en miembros inferiores, traccionar desde los hombros. Si los riñones no son el área de interés, el enfoque puede realizarse en PA estando el paciente de pie para reducir la dosis gonadal.
	
	\clearpage
	\onecolumn
	\begin{figure}
		\centering
		\includegraphics[width=0.7\linewidth]{imagenes/abd4}
		\caption{Desde sínfisis del pubis hacia arriba}
		\label{fig:abd3}
	\end{figure}
	
	\twocolumn
	%----------------------------------------------------------
	\clearpage
	
	\subsection{Abdomen AP en Bipedestación}
	
	\subsubsection*{1. Identificación}
	
	\textbf{Tipo:} Proyección básica. Forma par con la proyección AP en decúbito supino.
	
	\textbf{Indicaciones clínicas principales:}
	\begin{itemize}
		\item Búsqueda de aire libre subdiafragmático (neumoperitoneo)
		\item Obstrucción intestinal con evaluación de niveles hidroaéreos
		\item Abdomen agudo cuando el paciente puede mantenerse en pie
	\end{itemize}
	
	\subsubsection*{2. Técnica Base}
	
	\textbf{Receptor de imagen:} 35$\times$43 cm. Orientación \textbf{longitudinal} en todos los biotipos.
	
	\textbf{DFRI:} \textit{[dato no suministrado]}.
	
	\textbf{Técnica:} Intermedia. Igual justificación que en decúbito: el bajo contraste intrínseco abdominal requiere escala larga de grises con kVp intermedio. En paciente obeso se prefiere técnica básica por la mayor radiación dispersa generada por el gran espesor.
	
	\textbf{Factores técnicos:} kVp \textbf{[valor]} --- mAs \textbf{[valor]}. Medición del espesor a nivel de las crestas iliacas.
	
	\textbf{Bucky:} Bucky mural con grilla.
	
	\subsubsection*{3. Posicionamiento}
	
	\paragraph{Paciente}
	De pie, peso distribuido equitativamente entre ambos pies, brazos a los lados del cuerpo. Plano medio sagital perpendicular y centrado al receptor vertical. Ausencia de rotación de pelvis y hombros.
	
	\begin{figure}
		\centering
		\includegraphics[width=0.7\linewidth]{imagenes/abd}
		\caption{}
		\label{fig:abd}
	\end{figure}
	
	
	\paragraph{Región de interés}
	Abdomen y pelvis sin rotación: EIAS equidistantes al receptor. El colimador se abre en su totalidad a lo ancho para incluir toda la región.
	
	\paragraph{Rayo Central}
	Perpendicular al plano del receptor. Punto de entrada: \textbf{5 cm por encima de las crestas iliacas}, para abarcar las cúpulas diafragmáticas.
	
	\paragraph{Receptor de imagen}
	Centrado por encima del nivel de las crestas iliacas de forma que el borde superior incluya ambas cúpulas diafragmáticas y el borde inferior alcance la sínfisis del pubis. Si el paciente es muy alto, agregar una segunda proyección centrada en pelvis.
	
	\paragraph{Respiración}
	Al final de la espiración. El diafragma desciende, maximizando la visualización del espacio subdiafragmático y facilitando la detección de aire libre.
	
	\subsubsection*{4. Anatomía Visible}
	
	Cúpulas diafragmáticas bilaterales, espacio subdiafragmático, contorno hepático inferior, riñones, músculo psoas bilateral, marco intestinal con eventuales niveles hidroaéreos, columna lumbar y pelvis ósea hasta la sínfisis del pubis. La bipedestación es la posición de elección para toda búsqueda de aire libre; todo lo que sea buscar aire se hace de pie.
	
	\subsubsection*{5. Criterios de Evaluación}
	
	\textbf{Posicionamiento correcto:} Cúpulas diafragmáticas incluidas en el borde superior; sínfisis del pubis en el borde inferior; columna lumbar centrada en el receptor.
	
	\textbf{Ausencia de rotación:} Alas iliacas y agujeros obturadores simétricos; apófisis espinosas centradas en cuerpos vertebrales lumbares; caderas equidistantes a los bordes laterales.
	
	\textbf{Penetración y contraste:} Escala larga de grises; diferenciación entre gas, tejido blando, grasa y hueso; niveles hidroaéreos nítidos si están presentes.
	
	\subsubsection*{6. Puntos Críticos}
	
	\textbf{No incluir las cúpulas diafragmáticas:} Se pierde la posibilidad de detectar neumoperitoneo o patología subdiafragmática; corregir elevando el RI o recentrando el RC 5 cm por encima de las crestas iliacas.
	
	\textbf{No incluir la sínfisis del pubis:} La imagen carece de valor diagnóstico completo; si el paciente es muy alto, agregar proyección complementaria de pelvis.
	
	\textbf{Adaptaciones:} Si el paciente no puede mantenerse de pie, sustituir por decúbito lateral izquierdo con RC horizontal para evaluar aire libre.
	
	\textbf{Trucos:} Todo estudio que busque aire libre debe realizarse en bipedestación. La tracción desde tobillos (verificar ausencia de prótesis) mejora la alineación. Si los riñones no son área de interés, realizar en PA
	 para proteger las gónadas.
	 
\clearpage
\onecolumn

\begin{figure}
	\centering
	\includegraphics[width=0.7\linewidth]{imagenes/abd3}
	\caption{De cúpulas diafragmáticas hacia abajo}
	\label{fig:abd3}
\end{figure}


\clearpage
\twocolumn

%----------------------------------------------------------
%  Para agregar otro enfoque duplicá el bloque
%  desde \subsection hasta \textbf{Trucos}
%----------------------------------------------------------


%%%%%%%%%%%%%%%%%%%%%%%%%%%%%%%%%%%%%%%%%%%%%%%%%%%%%%%%%%%
%  SECCIÓN: Abdomen
%  Archivo: secciones/abdomen.tex
%  Sin preámbulo ni \begin{document} — lo maneja main.tex
	%%%%%%%%%%%%%%%%%%%%%%%%%%%%%%%%%%%%%%%%%%%%%%%%%%%%%%%%%%%
	
	\section{Abdomen}
	
	\subsection{AP Decúbito Lateral Izquierdo --- RC Horizontal (Frente Lateralizado)}
	
	\subsubsection*{1. Identificación}
	
	\textbf{Tipo:} Proyección complementaria. Suplanta la proyección frente en bipedestación cuando el paciente no puede ponerse de pie.
	
	\textbf{Indicaciones clínicas principales:}
	\begin{itemize}
		\item Sospecha de neumoperitoneo (aire libre intraperitoneal)
		\item Detección de niveles hidroaéreos en paciente sin tolerancia a la bipedestación
		\item Paciente en estado crítico o con movilidad muy reducida
	\end{itemize}
	
	\subsubsection*{2. Técnica Base}
	
	\textbf{Receptor de imagen:} 35$\times$43 cm. Orientación \textbf{transversal}, siguiendo el eje longitudinal de la región abdominal.
	
	\textbf{DFRI:} 100 cm. Distancia estándar para abdomen simple.
	
	\textbf{Técnica:} Alto kVp (técnica básica). Permite visualizar el aire libre como zona radiolúcida sobre el fondo gris de los órganos sólidos; el RC atraviesa el abdomen de forma horizontal para que el aire, al ascender por flotabilidad, quede proyectado en el margen superior de la cavidad.
	
	\textbf{Factores técnicos:} kVp \textbf{[dato no suministrado]} --- mAs \textbf{[dato no suministrado]}. Espesor medido a nivel de las crestas ilíacas.
	
	\textbf{Bucky:} Bucky de mesa con grilla.
	
	\subsubsection*{3. Posicionamiento}
	
	\paragraph{Paciente}
	Decúbito lateral izquierdo (perfil), apoyando el flanco izquierdo sobre la mesa o camilla. Se prefiere el lado izquierdo para evitar que el aire libre se confunda con la burbuja gástrica, que normalmente ocupa ese hemiabdomen. Si es posible, mantener al paciente en esta posición durante al menos 5 minutos antes de la exposición para permitir que el aire libre ascienda a su máxima altura y el líquido intraperitoneal se redistribuya hacia el flanco izquierdo.
	
	\begin{figure}[h]
		\centering
		\includegraphics[width=0.7\linewidth]{imagenes/abd5}
		\caption{}
		\label{fig:abd5}
	\end{figure}
	
	
	\paragraph{Región de interés}
	Crestas ilíacas centradas sobre el receptor de imagen. Rodillas parcialmente flexionadas, apoyadas una sobre la otra, para estabilizar al paciente. Brazos elevados por encima del nivel del diafragma y apoyados bajo la cabeza, de modo que no se superpongan sobre el contenido abdominal.
	
	\paragraph{Rayo Central}
	Horizontal, perpendicular al plano del receptor. Punto de entrada: \textbf{centro de la región abdominal, 5 cm por encima de las crestas ilíacas}.
	
	\paragraph{Receptor de imagen}
	Centrado a nivel de las crestas ilíacas. Debe incluir ambos diafragmas en el margen superior y la sínfisis del pubis en el margen inferior. El receptor se coloca vertical, paralelo al eje longitudinal del paciente.
	
	\paragraph{Respiración}
	Al final de la espiración. Reduce el volumen pulmonar y desciende los diafragmas, ampliando el espacio visible de la cavidad abdominal superior. Si el paciente está lúcido, solicitarle colaboración activa para sostener la apnea.
	
	\subsubsection*{4. Anatomía Visible}
	
	Cavidad abdominal completa con ambos diafragmas incluidos. Silueta hepática en el hemiabdomen derecho. Burbuja gástrica en el hemiabdomen izquierdo. Niveles hidroaéreos dentro de las asas intestinales. Aire libre intraperitoneal proyectado sobre la silueta del hígado (lado derecho, margen superior) cuando existe neumoperitoneo. El lado superior es el de mayor relevancia diagnóstica para detección de aire libre.
	
	\subsubsection*{5. Criterios de Evaluación}
	
	\textbf{Posicionamiento correcto:} Ambos diafragmas incluidos en el campo. RC horizontal confirmado por la visualización de niveles líquidos perfectamente horizontales dentro de las asas. Crestas ilíacas centradas en la imagen.
	
	\textbf{Ausencia de rotación:} Alas ilíacas simétricas. Bordes costales inferiores externos equidistantes respecto a la columna vertebral. Cuerpos vertebrales sin oblicuidad apreciable.
	
	\textbf{Penetración y contraste:} Escala de grises que permite distinguir el aire libre (negro) sobre el fondo gris de los órganos sólidos. Niveles hidroaéreos visibles tanto en posición superior (gas) como inferior (líquido) dentro de las asas intestinales. Aire intraperitoneal libre identificable sobre la silueta hepática.
	
	\subsubsection*{6. Puntos Críticos}
	
	\textbf{RC no horizontal:} Si el RC no es estrictamente horizontal, los niveles hidroaéreos se distorsionan y el aire libre puede no proyectarse en el margen correcto, llevando a falsos negativos. Verificar la horizontalidad del tubo antes de exponer.
	
	\textbf{Tiempo insuficiente en decúbito:} Si no se respetan los 5 minutos previos, el aire libre puede no haber ascendido lo suficiente y quedar oculto. Cuando la situación clínica lo permita, respetar el tiempo de espera.
	
	\textbf En pacientes con trauma abdominal severo o inestabilidad hemodinámica, la exposición se realiza con el mínimo movimiento posible; se puede posicionar el receptor sin mover al paciente del transporte si la situación lo requiere.
	
	\textbf{Trucos:} Visualizar siempre el lado superior de la imagen al evaluar: es donde se proyecta el aire libre. La burbuja gástrica en el lado izquierdo no debe confundirse con neumoperitoneo; el decúbito izquierdo aleja precisamente ese artefacto del área diagnóstica crítica (margen hepático derecho).
	
	\onecolumn
	
	\begin{figure}
		\centering
		\includegraphics[width=0.7\linewidth]{imagenes/abd6}
		\caption{}
		\label{fig:abd6}
	\end{figure}
	
	\twocolumn

\clearpage

	
	%----------------------------------------------------------
	%  Para agregar otro enfoque duplicá el bloque
	%  desde \subsection hasta \textbf{Trucos}
	%----------------------------------------------------------
	
	
%%%%%%%%%%%%%%%%%%%%%%%%%%%%%%%%%%%%%%%%%%%%%%%%%%%%%%%%%%%
%  SECCIÓN: Abdomen
%  Archivo: secciones/abdomen.tex
%  Sin preámbulo ni \begin{document} — lo maneja main.tex
	%%%%%%%%%%%%%%%%%%%%%%%%%%%%%%%%%%%%%%%%%%%%%%%%%%%%%%%%%%%
	
	\section{Abdomen}
	
	\subsection{Abdomen Perfil}
	
	\subsubsection*{1. Identificación}
	
	\textbf{Tipo:} Proyección complementaria. Forma par con la proyección frente de abdomen.
	
	\textbf{Indicaciones clínicas principales:}
	\begin{itemize}
		\item Localización tridimensional de cuerpo extraño
		\item Aneurisma o calcificación de aorta abdominal
		\item Neoplasias, ascitis y exploraciones con medio de contraste
	\end{itemize}
	
	\subsubsection*{2. Técnica Base}
	
	\textbf{Receptor de imagen:} 35 $\times$ 43 cm. Orientación \textbf{longitudinal} en relación con el eje mayor de la región.
	
	\textbf{DFRI:} 100 cm. Distancia estándar para estudios de abdomen.
	
	\textbf{Técnica:} Básica o contrastada según indicación. A mayor espesor, mayor kVp requerido.
	
	\textbf{Factores técnicos:} kVp \textbf{[dato no suministrado]} --- mAs \textbf{[dato no suministrado]}. Medir el espesor en la región de mayor grosor abdominal.
	
	\textbf{Bucky:} Bucky mural (bipedestación) o Bucky de mesa (decúbito lateral). Con grilla.
	
	\subsubsection*{3. Posicionamiento}
	
	\paragraph{Paciente}
	De pie (bipedestación) o en decúbito lateral derecho o izquierdo. Plano medio sagital paralelo al receptor. Codos flexionados con ambos brazos levantados por encima de la cabeza, alejándose de la región abdominal.
	
	\begin{figure}[h]
		\centering
		\includegraphics[width=0.7\linewidth]{imagenes/abd7}
		\caption{}
		\label{fig:abd7}
	\end{figure}
	
	
	\paragraph{Región de interés}
	Pelvis sin rotación. Rodillas flexionadas para reducir la lordosis lumbar; colocar apoyo bajo rodillas y tobillos en decúbito. Plano coronal del paciente perpendicular al receptor.
	
	\paragraph{Rayo Central}
	Perpendicular al plano del receptor. Punto de entrada: \textbf{5 cm por encima de las crestas ilíacas, al plano medio coronal}, de modo que incluya el diafragma.
	
	\paragraph{Receptor de imagen}
	Centrado a la altura de las crestas ilíacas (o 5 cm por encima). Debe incluir desde las cúpulas diafragmáticas hasta la pelvis, y los bordes anterior y posterior del abdomen.
	
	\paragraph{Respiración}
	Al final de la espiración. Permite mayor inclusión del diafragma y reduce la superposición de estructuras intestinales.
	
	\subsubsection*{4. Anatomía Visible}
	
	Asas intestinales, niveles hidroaéreos si presentes, aorta abdominal y sus calcificaciones, columna lumbar en perfil, cúpulas diafragmáticas, retroperitoneo y pelvis ósea. La columna debe verse recta y alineada con el receptor. Los huesos ilíacos deben aparecer superpuestos (con magnificación inevitable de uno respecto al otro por la divergencia del haz). No se incluye la sínfisis del pubis.
	
	\subsubsection*{5. Criterios de Evaluación}
	
	\textbf{Posicionamiento correcto:} Abdomen de perfil visible desde cúpulas diafragmáticas hasta pelvis; columna lumbar alineada con el receptor; sínfisis del pubis excluida.
	
	\textbf{Ausencia de rotación:} Bordes posteriores de los arcos costales superpuestos; bordes posteriores de las alas ilíacas superpuestos. Un ilíaco siempre presentará mayor magnificación que el otro por efecto geométrico.
	
	\textbf{Penetración y contraste:} Asas intestinales y niveles hidroaéreos diferenciables; tejidos blandos abdominales visibles en la región anterior y posterior.
	
	\subsubsection*{6. Puntos Críticos}
	
	\textbf{Rotación del paciente:} Produce asimetría de los arcos costales e ilíacos posteriores, impidiendo la localización tridimensional. Corregir alineando el plano coronal estrictamente perpendicular al receptor.
	
	\textbf{Centrado incorrecto del RC:} Si el RC se dirige a la altura exacta de crestas sin compensar, el diafragma puede quedar excluido. Subir 5 cm por encima de las crestas para asegurar su inclusión.
	
	\textbf{Adaptaciones:} En paciente politraumatizado que no puede moverse se utiliza el perfil tangencial (ver siguiente subsección). No se emplea esta proyección para diagnóstico de oclusión intestinal; en ese caso usar frente en decúbito dorsal más frente de pie.
	
	\textbf{Trucos:} Usar el cátodo en la parte superior del abdomen (región más gruesa) para obtener exposición uniforme aprovechando el efecto anódico. Colocar apoyo bajo rodillas y tobillos en decúbito lateral para mantener el plano coronal perpendicular y evitar rotación involuntaria.
	
	\clearpage
	\onecolumn
	
	\begin{figure}
		\centering
		\includegraphics[width=0.7\linewidth]{imagenes/abd8}
		\caption{}
		\label{fig:abd8}
	\end{figure}
	
	
	
	\twocolumn
	\clearpage	
	%----------------------------------------------------------
	
	\subsection{Abdomen Perfil Tangencial}
	
	\subsubsection*{1. Identificación}
	
	\textbf{Tipo:} Proyección complementaria. Variante del perfil de abdomen para pacientes que no pueden adoptar bipedestación ni movilizarse.
	
	\textbf{Indicaciones clínicas principales:}
	\begin{itemize}
		\item Paciente que no puede ponerse de pie
		\item Paciente politraumatizado con imposibilidad de movimiento
		\item Sospecha de neumoperitoneo en paciente inmovilizado
	\end{itemize}
	
	\subsubsection*{2. Técnica Base}
	
	\textbf{Receptor de imagen:} 35 $\times$ 43 cm. Orientación \textbf{longitudinal} en relación con el eje mayor de la región.
	
	\textbf{DFRI:} 100 cm. Distancia estándar para estudios de abdomen.
	
	\textbf{Técnica:} Básica o contrastada según indicación.
	
	\textbf{Factores técnicos:} kVp \textbf{[dato no suministrado]} --- mAs \textbf{[dato no suministrado]}. Medir el espesor en la región de mayor grosor abdominal.
	
	\textbf{Bucky:} Con grilla. Adaptado a la posición del paciente (Bucky mural con tubo horizontal).
	
	\subsubsection*{3. Posicionamiento}
	
	\paragraph{Paciente}
	\textit Paciente decúbito supino, con brazos fuera del campo visual
	
	\begin{figure}[h]
		\centering
		\includegraphics[width=0.7\linewidth]{imagenes/abd9}
		\caption{}
		\label{fig:abd9}
	\end{figure}
	
	\paragraph{Rayo Central}
	Perpendicular al plano del receptor, dirigido horizontalmente. Punto de entrada: \textbf{a la altura de las crestas ilíacas, al plano medio coronal}.
	
	\paragraph{Receptor de imagen}
	Centrado a la altura de las crestas ilíacas. Debe incluir desde las cúpulas diafragmáticas hasta la pelvis y los bordes anterior y posterior del abdomen.
	
	\paragraph{Respiración}
	\textit{suspendida en al final de una experiacion}
	
	\subsubsection*{4. Anatomía Visible}
	
	Abdomen de perfil desde cúpulas diafragmáticas hasta pelvis. Huesos ilíacos superpuestos como indicador de ausencia de rotación. En caso de neumoperitoneo, el aire libre se visualiza en la región anterior del abdomen (entre la pared abdominal anterior y el hígado/vísceras).
	
	\subsubsection*{5. Criterios de Evaluación}
	
	\textbf{Posicionamiento correcto:} Abdomen de perfil visible desde cúpulas diafragmáticas hasta pelvis.
	
	\textbf{Ausencia de rotación:} Huesos ilíacos superpuestos.
	
	\textbf{Penetración y contraste:} Diferenciación de tejidos blandos; región anterior del abdomen visible para detectar aire libre subdiafragmático anterior.
	
	\subsubsection*{6. Puntos Críticos}
	
	\textbf{Indicación incorrecta:} El perfil tangencial no reemplaza al perfil estándar en pacientes móviles; se reserva exclusivamente para politraumatizados o pacientes que no pueden moverse.
	
	\textbf{Rotación del paciente:} Impide la superposición de los ilíacos y dificulta la identificación de neumoperitoneo. Verificar alineación antes de la exposición.
	
	\textbf{Adaptaciones:} Técnica diseñada específicamente para pacientes con movilidad reducida o nula. El tubo se orienta horizontalmente con el receptor vertical junto al paciente.
	
	\textbf{Trucos:} La presencia de aire libre anterior al abdomen en esta proyección es indicador sensible de neumoperitoneo en el paciente inmovilizado.
	
		\clearpage
	\onecolumn
	
	\begin{figure}
		\centering
		\includegraphics[width=0.7\linewidth]{imagenes/abd10}
		\caption{}
		\label{fig:abd8}
	\end{figure}
	
	
	
	\twocolumn
	\clearpage	
	%----------
	
	%----------------------------------------------------------
	%  Para agregar otro enfoque duplicá el bloque
	%  desde \subsection hasta \textbf{Trucos}
	%----------------------------------------------------------
	
	
%%%%%%%%%%%%%%%%%%%%%%%%%%%%%%%%%%%%%%%%%%%%%%%%%%%%%%%%%%%
%  SECCIÓN: Abdomen — Paciente Pediátrico
%  Archivo: secciones/abdomen_pediatrico.tex
%  Sin preámbulo ni \begin{document} — lo maneja main.tex
	%%%%%%%%%%%%%%%%%%%%%%%%%%%%%%%%%%%%%%%%%%%%%%%%%%%%%%%%%%%
	
	\section{Abdomen — Paciente Pediátrico}
	
	\subsection*{Consideraciones Clínicas}
	
	\textbf{Selección del estudio según diagnóstico clínico (DC):}
	El DC condiciona tanto la proyección como la posición del paciente. Dos situaciones frecuentes requieren criterio específico:
	
	\begin{itemize}
		\item \textbf{Abdomen agudo:} proyección de pie, incluyendo desde la sínfisis púbica hasta el diafragma.
		\item \textbf{Cuerpo extraño (CE):} frente y perfil de pie. Puede requerir campo desde peñascos hasta sínfisis púbica según tiempo de ingesta.
	\end{itemize}
	
	\textbf{Paciente que no puede estar de pie:} se realiza sentado o en decúbito lateral con RC horizontal, siempre que se requiera evaluar niveles hidroaéreos.
	
	\textbf{CE en decúbito supino:} aceptable en frente y perfil cuando no es posible la bipedestación ni la sedestación.
	
	\textbf{Protección gonadal:} aplicar siempre que no interfiera con el área de interés. En pacientes pediátricos se prioriza obtener una exposición diagnóstica correcta sin protección antes que intentar colocarla y comprometer la región de interés.
	
	\textbf{Campo de imagen:} toda la región de interés debe entrar en un único chasis.
	
	\medskip
	
	\textbf{¿Por qué de pie o sentado en CE?} Permite evaluar si el objeto puede haber perforado una víscera y determinar su localización exacta mediante niveles.
	
	\textbf{¿Por qué incluir desde peñascos hacia abajo?} Si la ingesta es reciente (menos de 2 horas), el CE puede no haber descendido al tracto digestivo y encontrarse aún en cuello o región torácica superior.
	
	\medskip
	
	\textbf{Interrogatorio previo:} preguntar a los padres o al niño (si puede comunicarse) el tiempo transcurrido desde la ingesta. Este dato orienta el nivel de búsqueda:
	
	\begin{itemize}
		\item \textbf{Menos de 2 hs:} cuello y región comprendida entre faringe y carina (vía aérea superior y esófago proximal). Se pide Rx de cuello, tórax y abdomen.
		\item \textbf{Primeras 12 hs:} estómago (tórax con porción superior de abdomen).
		\item \textbf{Más de 12 hs:} intestino delgado.
		\item \textbf{Entre 12 y 24 hs:} intestino grueso.
		\item \textbf{Más de 24 hs:} si no se visualiza, probable expulsión o localización en recto inminente.
	\end{itemize}
	
	%----------------------------------------------------------
	%  Continuar con \subsection{} para cada enfoque
	%----------------------------------------------------------
	
\clearpage
%%%%%%%%%%%%%%%%%%%%%%%%%%%%%%%%%%%%%%%%%%%%%%%%%%%%%%%%%%%
%  SECCIÓN: Aparato Urinario
%  Archivo: secciones/aparato_urinario.tex
%  Sin preámbulo ni \begin{document} — lo maneja main.tex
	%%%%%%%%%%%%%%%%%%%%%%%%%%%%%%%%%%%%%%%%%%%%%%%%%%%%%%%%%%%
	
	\section{Abdomen}
	
	\subsection{Aparato Urinario Simple (AUS) --- Proyección AP}
	
	\subsubsection*{1. Identificación}
	
	\textbf{Tipo:} Proyección básica. Radiografía exploratoria sin contraste, equivalente a un frente de abdomen. Constituye la proyección preliminar obligatoria antes de estudios con medio de contraste (UIV, uretrocistografía, etc.).
	
	\textbf{Indicaciones clínicas principales:}
	\begin{itemize}
		\item Litiasis renal o ureteral: detección de cálculos por su densidad cálcica pura (el contraste los encubriría)
		\item Valoración de calcificaciones anómalas en la vía urinaria
		\item Estudio preliminar previo a la administración de medio de contraste para identificar signos basales (obstrucción, hidronefrosis, tumores, infecciones)
	\end{itemize}
	
	\subsubsection*{2. Técnica Base}
	
	\textbf{Receptor de imagen:} 35$\times$43 cm. Orientación \textbf{longitudinal} para incluir desde T12 hasta la sínfisis del pubis. En pacientes muy altos, se realiza una primera exposición con este chasis y se complementa con un segundo receptor de 24$\times$30 cm en orientación \textbf{transversal}, centrado 5--7 cm por encima del borde superior de la sínfisis del pubis, para incluir la vejiga.
	
	\textbf{DFRI:} 100 cm. Distancia estándar para radiografía de abdomen.
	
	\textbf{Técnica:} Técnica intermedia. Las densidades del abdomen son muy homogéneas (bajo contraste intrínseco), por lo que se requiere una escala larga de grises para diferenciar el contorno de las vísceras (riñones), tejido adiposo, aire y estructuras óseas.
	
	\textbf{Factores técnicos:} Medición del espesor a nivel de las crestas ilíacas (L4).
	
	\textbf{Bucky:} Bucky de mesa con grilla.
	
	\subsubsection*{3. Posicionamiento}
	
	\paragraph{Paciente}
	Decúbito supino sobre la mesa de examen. El plano medio sagital (PMS) perpendicular a la mesa y centrado sobre la línea media del receptor. Los brazos colocados de manera que no proyecten sombras sobre el receptor de imagen.
	
	\begin{figure}[h]
		\centering
		\includegraphics[width=0.7\linewidth]{imagenes/aus1}
		\caption{}
		\label{fig:aus1}
	\end{figure}
	
	
	\paragraph{Región de interés}
	Sin rotación de pelvis ni tronco: las espinas ilíacas anterosuperiores (EIAS) deben quedar equidistantes a la mesa. Leve flexión de rodillas con apoyo para aliviar la tensión lumbar. La región a incluir comprende desde T12 (polo superior de los riñones) hasta la sínfisis del pubis (incluida la vejiga, que puede extenderse por debajo de la sínfisis cuando está llena).
	
	\paragraph{Rayo Central}
	Perpendicular al plano del receptor. Punto de entrada: \textbf{a la altura de las crestas ilíacas (L4), sobre el plano medio coronal}.
	
	\paragraph{Receptor de imagen}
	Centrado a la altura de las crestas ilíacas. Borde superior en T12 para incluir los polos superiores renales; borde inferior en la sínfisis del pubis para incluir vejiga completa. En el enfoque complementario de vejiga: receptor de 24$\times$30 cm centrado 5--7 cm por encima del borde superior de la sínfisis del pubis.
	
	\paragraph{Respiración}
	Espiración máxima, suspendida. La espiración desciende el diafragma y reduce la superposición de estructuras torácicas basales sobre los riñones; la apnea elimina el movimiento de las vísceras abdominales durante la exposición.
	
	\subsubsection*{4. Anatomía Visible}
	
	Silueta completa de ambos riñones (polos superiores a nivel de T12, inferiores aproximadamente en L3--L4). Columna vertebral lumbar centrada en la imagen. Contornos de los músculos psoas. Sínfisis del pubis y vejiga urinaria. Alas ilíacas simétricas. En presencia de litiasis: opacidades cálcicas a lo largo del trayecto ureteral (desde pelvis renal hasta vejiga) o en el parénquima renal.
	
	\subsubsection*{5. Criterios de Evaluación}
	
	\textbf{Posicionamiento correcto:} Columna vertebral centrada en la radiografía. Inclusión completa desde T12 hasta la sínfisis del pubis. Silueta renal bilateral visible. Vejiga incluida en el campo.
	
	\textbf{Ausencia de rotación:} Alas ilíacas simétricas. EIAS equidistantes al eje vertical de la imagen. Apófisis espinosas centradas respecto a los pedículos vertebrales.
	
	\textbf{Penetración y contraste:} Escala larga de grises que permita diferenciar el contorno renal del tejido adiposo perirrenal, el gas intestinal y las estructuras óseas. Ausencia de artefactos por ropa interior u objetos metálicos del paciente.
	
	\subsubsection*{6. Puntos Críticos}
	
	\textbf{Paciente demasiado alto para un solo chasis:} La vejiga y la sínfisis del pubis quedan fuera del campo. Corrección: realizar primero la exposición de abdomen completo y luego un enfoque complementario de vejiga con receptor de 24$\times$30 cm transversal, centrado 5--7 cm por encima del borde superior de la sínfisis del pubis.
	
	\textbf{Rotación de pelvis o tronco:} Proyección asimétrica de las alas ilíacas y falsa lateralización de estructuras urinarias, dificultando la localización de cálculos. Corrección: verificar que las EIAS sean equidistantes a la mesa antes de disparar.
	
	\textbf{Adaptaciones:} En paciente con dificultad para el decúbito supino, puede realizarse en bipedestación (AP de pie); en ese caso verificar que no haya desplazamiento renal por efecto gravitacional que altere la región incluida.
	
	\textbf{Trucos:} Solicitar al paciente que retire ropa interior con elementos metálicos (hebillas, ganchos) para evitar artefactos sobre la región pelviana. Recordar que los cálculos NO se ven con medio de contraste: esta proyección es la única que los muestra por su densidad cálcica pura. No es indispensable incluir las cúpulas diafragmáticas; la región crítica es el polo inferior, la vejiga y la sínfisis.
	
	\begin{figure}[h]
		\centering
		\includegraphics[width=0.7\linewidth]{imagenes/aus3}
		\caption{}
		\label{fig:aus3}
	\end{figure}
	
	\textbf{Caso:} Pct. rotado, afecta el valor diagnóstico (volver a repetirse, porque se evalúa morfología y función, si un lado está más inclinado es posible que el MC llegue mas rapido por eso mismo, falseando así el resultado).
	
	
	\clearpage
	\onecolumn
	
	\begin{figure}
		\centering
		\includegraphics[width=0.7\linewidth]{imagenes/aus2}
		\caption{}
		\label{fig:aus2}
	\end{figure}
	
	
	
	\clearpage
	\twocolumn
	%----------------------------------------------------------
	%  Para agregar otro enfoque duplicá el bloque
	%  desde \subsection hasta \textbf{Trucos}
	%----------------------------------------------------------
	
	
%%%%%%%%%%%%%%%%%%%%%%%%%%%%%%%%%%%%%%%%%%%%%%%%%%%%%%%%%%%
%  SECCIÓN: Vejiga Urinaria
%  Archivo: secciones/vejiga.tex
%  Sin preámbulo ni \begin{document} — lo maneja main.tex
	%%%%%%%%%%%%%%%%%%%%%%%%%%%%%%%%%%%%%%%%%%%%%%%%%%%%%%%%%%%
	
	\section{Abdomen}
	
	\subsection{Vejiga Urinaria - Proyección AP Axial}
	
	\subsubsection*{1. Identificación}
	
	\textbf{Tipo:} Proyección complementaria. Se realiza como complemento cuando el chasis 35$\times$43 no alcanza para incluir la vejiga, o para valorar específicamente la vejiga.
	
	\textbf{Indicaciones clínicas principales:}
	\begin{itemize}
		\item Divertículo vesical
		\item Reflujo vesicoureteral
		\item Evaluación de próstata en hombres
	\end{itemize}
	
	\subsubsection*{2. Técnica Base}
	
	\textbf{Receptor de imagen:} 24$\times$30 cm. Orientación \textbf{transversal} en todos los biotipos.
	
	\textbf{DFRI:} 100 cm. Distancia estándar para estudios de pelvis con técnica convencional.
	
	\textbf{Técnica:} Intermedia. Permite diferenciar la densidad de la vejiga llena con medio de contraste respecto a las estructuras óseas (pelvis y sacro) que se encuentran por detrás.

	
	\textbf{Bucky:} Bucky de mesa con grilla.
	
	\subsubsection*{3. Posicionamiento}
	
	\paragraph{Paciente}
	Decúbito supino sobre la mesa (también puede realizarse en bipedestación). Plano medio sagital perpendicular al receptor. Brazos ubicados de modo que no proyecten sombra sobre el receptor de imagen.
	
	\begin{figure}[h]
		\centering
		\includegraphics[width=0.7\linewidth]{imagenes/aus4}
		\caption{}
		\label{fig:aus4}
	\end{figure}
	
	
	\paragraph{Región de interés}
	Pelvis sin rotación: las espinas ilíacas anterosuperiors (EIAS) equidistantes al receptor. Piernas extendidas para lograr que la zona lumbosacra se arquee e incline hacia abajo los huesos pélvicos anteriores. Hombros en un mismo plano.
	
	\paragraph{Rayo Central}
	Angulado 10--15° en sentido caudal respecto al plano del receptor. La angulación exacta depende del grado de lordosis lumbar: a mayor lordosis, menor angulación requerida. Punto de entrada: \textbf{5 cm por encima del borde superior de la sínfisis del pubis}.
	
	\paragraph{Receptor de imagen}
	Centrado 5 cm por encima del borde superior de la sínfisis del pubis (referencia de altura equivalente al trocánter mayor del fémur). Debe incluir los extremos distales de los uréteres, la vejiga y la parte proximal de la uretra.
	
	\paragraph{Respiración}
	Espiración máxima, luego suspendida. Reduce el contenido abdominal sobre la región de interés y mejora la visualización de la vejiga.
	
	\subsubsection*{4. Anatomía Visible}
	
	Extremos distales de los uréteres (unión ureterovesical), vejiga urinaria con medio de contraste, y parte proximal de la uretra. Los huesos del pubis deben proyectarse por debajo del cuello vesical y la uretra proximal, sin superponerse a estas estructuras. La angulación caudal del RC desplaza la sínfisis del pubis hacia caudal, dejando libre la región del cuello vesical.
	
	\subsubsection*{5. Criterios de Evaluación}
	
	\textbf{Posicionamiento correcto:} La sínfisis del pubis aparece proyectada por debajo de la vejiga, sin superponer el cuello vesical ni la zona de los uréteres distales. El medio de contraste se visualiza claramente en vejiga, uréteres distales y uretra proximal.
	
	\textbf{Ausencia de rotación:} Las EIAS equidistantes respecto a la línea media; pelvis simétrica en la imagen.
	
	\textbf{Penetración y contraste:} Escala de contraste corta (alta densidad relativa) que permita diferenciar claramente el medio de contraste en la vejiga respecto a las estructuras óseas de la pelvis y el sacro subyacente.
	
	\subsubsection*{6. Puntos Críticos}
	
	\textbf{Angulación insuficiente del RC:} La sínfisis del pubis se superpone al cuello vesical y a los uréteres distales, impidiendo su evaluación. Aumentar la angulación caudal dentro del rango 10--15°.
	
	\textbf{Piernas no extendidas:} La zona lumbosacra no se arquea correctamente y los huesos pélvicos anteriores no se inclinan hacia abajo, reduciendo la efectividad de la angulación. Verificar que las piernas estén en extensión completa antes de exponer.
	
	\textbf{Adaptaciones:} En pacientes con mayor lordosis lumbar puede requerirse menor angulación (incluso por debajo de 10°). En bipedestación, verificar que el plano medio sagital se mantenga perpendicular al receptor.
	
	\textbf{Trucos:} El trocánter mayor del fémur sirve como referencia de altura para centrar el RI cuando no se palpa con facilidad el borde superior de la sínfisis. La angulación caudal es el elemento diferenciador de esta proyección respecto a la AP estándar de pelvis.
	
	
	\clearpage
	\onecolumn
	\begin{figure}
		\centering
		\includegraphics[width=0.7\linewidth]{imagenes/aus5}
		\caption{}
		\label{fig:aus5}
	\end{figure}
	
	\twocolumn
	\clearpage
	
	
	
	%----------------------------------------------------------
	%  Para agregar otro enfoque duplicá el bloque
	%  desde \subsection hasta \textbf{Trucos}
	%----------------------------------------------------------

%%%%%%%%%%%%%%%%%%%%%%%%%%%%%%%%%%%%%%%%%%%%%%%%%%%%%%%%%%%
%  SECCIÓN: Vejiga Urinaria
%  Archivo: secciones/vejiga.tex
%  Sin preámbulo ni \begin{document} — lo maneja main.tex
	%%%%%%%%%%%%%%%%%%%%%%%%%%%%%%%%%%%%%%%%%%%%%%%%%%%%%%%%%%%
	
	\section{Abdomen}
	
	\subsection{Vejiga - OPI/OPD}
	
	\subsubsection*{1. Identificación}
	
	\textbf{Tipo:} Proyección complementaria. Se realizan ambas oblicuas (OPD y OPI) como estudio comparativo bilateral.
	
	\textbf{Indicaciones clínicas principales:}
	\begin{itemize}
		\item Valoración de la unión ureterovesical (pasaje ureterovesical)
		\item Estudio de uréteres distales sin superposición con estructuras pélvicas
		\item Solicitud en el contexto de urografía excretora
	\end{itemize}
	
	\subsubsection*{2. Técnica Base}
	
	\textbf{Receptor de imagen:} 24$\times$30 cm. Orientación \textbf{longitudinal} (vertical).
	
	\textbf{DFRI:} 100 cm. Distancia estándar para proyecciones de abdomen y pelvis.
	
	\textbf{Técnica:} Intermedia. La rotación del paciente aumenta el espesor efectivo atravesado por el haz; se requiere mayor kVp para diferenciar la morfología interna de la vejiga sin que el medio de contraste sature la imagen.
	
	
	\textbf{Bucky:} Bucky de mesa con grilla.
	
	\subsubsection*{3. Posicionamiento}
	
	\paragraph{Paciente}
	Decúbito supino sobre la mesa de examen.
	
	\begin{figure}[h]
		\centering
		\includegraphics[width=0.7\linewidth]{imagenes/aus6}
		\caption{}
		\label{fig:aus6}
	\end{figure}
	
	
	\paragraph{Región de interés}
	Se gira al paciente 40--60° hacia el lado de interés (OPD: lado derecho apoyado; OPI: lado izquierdo apoyado). Para visualización óptima del pasaje ureterovesical, la rotación se reduce a 35°, ángulo que corresponde a la llegada del uréter a la vejiga respecto al plano coronal. La parte superior del muslo ipsilateral se extiende y abduce para evitar superposición sobre la zona vesical.
	
	\paragraph{Rayo Central}
	Perpendicular al plano del receptor. Punto de entrada: \textbf{5 cm por encima del borde superior de la sínfisis del pubis y 5 cm medial a la EIAS ipsilateral}. En estudios miccionales: perpendicular a la altura de la sínfisis púbica.
	
	\paragraph{Receptor de imagen}
	Centrado 5 cm por encima del borde superior de la sínfisis púbica y aproximadamente 5 cm medial a la EIAS (referencia Merrill: a la altura de las crestas ilíacas). Incluir los extremos distales de los uréteres, la vejiga completa y la parte proximal de la uretra.
	
	\paragraph{Respiración}
	Suspendida al final de la espiración. Estabiliza los órganos pélvicos durante la exposición.
	
	\subsubsection*{4. Anatomía Visible}
	
	Extremo distal de los uréteres, vejiga urinaria y parte proximal de la uretra. La unión ureterovesical del lado apoyado se proyecta libre de superposición. Los huesos del pubis deben proyectarse por debajo del cuello vesical y la uretra proximal. La parte superior del muslo no debe superponerse sobre la zona vesical.
	
	\subsubsection*{5. Criterios de Evaluación}
	
	\textbf{Posicionamiento correcto:} Huesos del pubis proyectados por debajo del cuello vesical y la uretra proximal. Extremos distales de los uréteres, vejiga y uretra proximal incluidos en el campo.
	
	\textbf{Ausencia de rotación:} El agujero obturador del lado apoyado aparece cerrado o estrecho; el del lado elevado aparece abierto. El ala ilíaca del lado apoyado se ve desplegada.
	
	\textbf{Penetración y contraste:} Escala de contraste semi-corta que muestre con claridad el medio de contraste en la vejiga, los uréteres distales y la uretra proximal, sin saturación de la imagen.
	
	\subsubsection*{6. Puntos Críticos}
	
	\textbf{Rotación incorrecta:} Menos de 35° superpone el uréter distal con la vejiga e impide ver el pasaje ureterovesical. Más de 60° pierde la referencia anatómica del lado de interés. Usar 35° específicamente para el pasaje ureterovesical; 40--60° para evaluación general.
	
	\textbf{Superposición del muslo:} Si no se extiende y abduce la parte superior del muslo, se cubre la región vesical inferior y la uretra proximal. Corregir la posición antes de la exposición.
	
	\textbf{Adaptaciones:} En pacientes con limitación de movilidad, usar cuñas radiolúcidas bajo el lado a elevar. En trauma pélvico, no forzar la rotación; documentar el grado alcanzado.
	
	\textbf{Trucos:} Para identificar el lado oblicuado en la radiografía: agujero obturador cerrado indica el lado hacia el que está oblicuado el paciente; agujero obturador abierto indica el lado elevado.
	
	\clearpage
	\onecolumn
	
	\begin{figure}
		\centering
		\includegraphics[width=0.7\linewidth]{imagenes/aus7}
		\caption{}
		\label{fig:aus7}
	\end{figure}
	
	\clearpage
	\twocolumn
	
	%----------------------------------------------------------
	%  Para agregar otro enfoque duplicá el bloque
	%  desde \subsection hasta \textbf{Trucos}
	%----------------------------------------------------------


%%%%%%%%%%%%%%%%%%%%%%%%%%%%%%%%%%%%%%%%%%%%%%%%%%%%%%%%%%%
%  SECCIÓN: Abdomen / Urografía
%  Archivo: secciones/abdomen.tex
%  Sin preámbulo ni \begin{document} — lo maneja main.tex
	%%%%%%%%%%%%%%%%%%%%%%%%%%%%%%%%%%%%%%%%%%%%%%%%%%%%%%%%%%%
	
	\section{Abdomen}
	
	\subsection{Urografía Intravenosa (Excretora) - OPD / OPI}
	
	\subsubsection*{1. Identificación}
	
	\textbf{Tipo:} Proyección complementaria. Forma par con la proyección AP de frente de la UIV.
	
	\textbf{Indicaciones clínicas principales:}
	\begin{itemize}
		\item Obstrucción del riñón elevado
		\item Signos de infección del tracto urinario
		\item Traumatismos en la cara inferior del uréter
	\end{itemize}
	
	\subsubsection*{2. Técnica Base}
	
	\textbf{Receptor de imagen:} 35$\times$43 cm. Orientación \textbf{longitudinal}.
	
	\textbf{DFRI:} 100 cm. Distancia estándar para estudios abdominales.
	
	\textbf{Técnica:} Intermedia o básica. \textit{[dato no suministrado sobre efecto en imagen]}.
	
	
	\textbf{Bucky:} Bucky de mesa con grilla.
	
	\subsubsection*{3. Posicionamiento}
	
	\paragraph{Paciente}
	Decúbito supino, girado de manera parcial hacia el lado derecho o izquierdo según el lado de interés. El giro permite separar el riñón elevado de la columna vertebral y proyectarlo en perfil.
	
	\begin{figure}[h]
		\centering
		\includegraphics[width=0.7\linewidth]{imagenes/aus8}
		\caption{}
		\label{fig:aus8}
	\end{figure}
	
	
	\paragraph{Región de interés}
	Rotar el cuerpo 30° para obtener una oblicua posterior derecha (OPD) o izquierda (OPI). Flexionar la pierna del lado elevado para estabilizar el cuerpo. Elevar el brazo del lado elevado y cruzarlo sobre el tórax. Centrar la columna vertebral sobre la línea media del receptor.
	
	\paragraph{Rayo Central}
	Perpendicular al plano del receptor. Punto de entrada: \textbf{5 cm lateral a la línea media del paciente, hacia el lado elevado, a la altura de las crestas ilíacas}.
	
	\paragraph{Receptor de imagen}
	Centrado de modo que el campo incluya desde los riñones en el borde superior hasta la sínfisis del pubis en el borde inferior.
	
	\paragraph{Respiración}
	Suspendida en espiración. La espiración completa desciende el diafragma y reduce la superposición de estructuras abdominales sobre los riñones.
	
	\subsubsection*{4. Anatomía Visible}
	
	Lado elevado (alejado del RI): riñón paralelo al receptor, proyectado de frente sin superposición con la columna vertebral; uréter distal visible. Lado apoyado (próximo al RI): riñón en perfil; uréter proximal y medio proyectados lejos de la columna, despejados. Sínfisis del pubis visible en el borde inferior del campo.
	
	\subsubsection*{5. Criterios de Evaluación}
	
	\textbf{Posicionamiento correcto:} Riñón elevado paralelo al plano del receptor y separado de la columna; riñón apoyado en perfil; sínfisis del pubis incluida en el borde inferior; ambos riñones dentro del campo superior.
	
	\textbf{Ausencia de rotación:} El grado de oblicuidad es correcto cuando el riñón elevado se muestra completamente de frente sin superponerse a la columna y el riñón apoyado aparece en perfil neto.
	
	\textbf{Penetración y contraste:} \textit{[dato no suministrado]}.
	
	\subsubsection*{6. Puntos Críticos}
	
	\textbf{Oblicuidad insuficiente:} El riñón elevado se superpone a la columna vertebral, impidiendo la evaluación del uréter distal; corregir aumentando la rotación hasta 30°.
	
	\textbf{Centrado incorrecto del RC:} Desplazar el punto de entrada 5 cm hacia el lado elevado es crítico; un RC centrado en la línea media distorsiona la proyección del riñón de interés.
	
	\textbf{Adaptaciones:} En pacientes con limitación de movilidad o trauma, apoyar la oblicua con cuñas radiolúcidas para mantener los 30° de rotación sin esfuerzo del paciente.
	
	\textbf{Trucos:} Verificar antes de exponer que la columna esté centrada sobre la línea media del RI; esto asegura que el riñón elevado quede correctamente lateralizado respecto al RC.
	
	\clearpage
	\onecolumn
	
	\begin{figure}
		\centering
		\includegraphics[width=0.7\linewidth]{imagenes/aus9}
		\caption{}
		\label{fig:aus9}
	\end{figure}
	
	\clearpage
	\twocolumn
	
	
	%----------------------------------------------------------
	%  Para agregar otro enfoque duplicá el bloque
	%  desde \subsection hasta \textbf{Trucos}
	%----------------------------------------------------------